\section{Schlussfolgerung}

Das Projekt zeigt, dass ein auf KITTI vortrainiertes 3D-Detektionsmodell durch Finetuning erfolgreich auf eine fest installierte Tiefgaragenszene angepasst werden kann. Bezogen auf die vier eingangs gestellten Forschungsfragen ergibt sich folgendes Bild.

\paragraph{Wirkung des Finetunings.} Ausschließlich vortrainierte Modelle bleiben auf diesen Aufnahmen bei einem mAP von höchstens 0{,}058, das hier finetunete Modell erreicht 0{,}798 -- ein Abstand von 0{,}740 mAP. Der Abstand zwischen den beiden untersuchten Architekturen beträgt demgegenüber 0{,}047 bis 0{,}118 mAP und ist damit sechs- bis sechzehnmal kleiner. Die beiden Zahlen sind unterschiedlich abgesichert. Der Architekturvergleich ist kontrolliert: Beide Modelle unterscheiden sich ausschließlich in der Detektorarchitektur und stimmen in allen übrigen Festlegungen überein; seine Grenze ist die Verallgemeinerung, da nur zwei Vertreter derselben gitterbasierten Familie auf einer einzigen Szene geprüft wurden. Der Baselinewert dagegen stammt aus einer Vorgängerarbeit, die nicht unter diesem Protokoll und möglicherweise nicht mit derselben Datenkonvertierung und Metrik erhoben wurde; der kausale Beitrag des Finetunings ist damit nicht kontrolliert bestimmt. Bei einem Unterschied dieser Größe ist die Richtung dennoch eindeutig, und für vergleichbare stationäre Aufbauten legt das Ergebnis nahe, den Aufwand zunächst in Zieldaten und deren Annotation zu investieren, bevor Architekturen verglichen werden.

\paragraph{Einzelsensoren gegenüber Fusion.} PointPillars liefert mit der fusionierten Punktwolke in jedem Cross-Validation-Fold den höchsten Gesamt-mAP. Über alle Objekte beträgt der Vorsprung im Fold-Mittel 0{,}021. Beschränkt auf die bewegten Objekte und auf derselben Aggregationsstufe berechnet, wächst er auf 0{,}040 in der offiziellen und 0{,}069 in der labelbasierten Auswertung. Im Gesamtwert fällt der Unterschied zwischen den Ansichten damit kleiner aus. Dieser Wert bewertet zugleich die Wiedererkennung zahlreicher wiederkehrender geparkter Fahrzeuge; wie stark diese den Abstand beeinflussen, lässt sich wegen der nichtlinearen AP-Berechnung nicht isoliert bestimmen.

\paragraph{Finden gegenüber Lokalisieren.} Der Fusionsvorteil liegt in der räumlichen Präzision, nicht im Auffinden. Bei AP\textsubscript{30} erreichen alle Ansichten mindestens 0{,}65 und überwiegend mehr als 0{,}80; jede bewegte Klasse wird grundsätzlich aus jeder Perspektive gefunden. Auf den einzelnen Testexperimenten ist bei AP\textsubscript{30} sogar \texttt{os0} die stärkste und gleichmäßigste Ansicht (Mittel 0{,}914 bei einer Standardabweichung von 0{,}023 gegenüber 0{,}871 bei 0{,}063 für \texttt{merged}). Erst bei AP\textsubscript{60} führt \texttt{merged}, bei Person und Fahrrad mit erheblichem Abstand und in beiden Auswertungsvarianten. Beim Fahrrad sind die Varianten sogar identisch, weil keine statischen Fahrräder annotiert sind. Nur beim Auto kehrt die offizielle Variante die Rangfolge um, was auf ein aufgeklärtes Zuordnungsartefakt an einem geparkten Fahrzeug zurückgeht.

\paragraph{Fehlerbilder und Laufzeitgrenzen.} Die dominierenden Fehler sind Geisterboxen an Positionen, an denen in den Trainingsaufnahmen Objekte vorkamen, sowie klassenübergreifende Mehrfachhypothesen an dünn abgetasteten Fahrzeugkanten. Beide verweisen auf die Ein-Szenen-Datenbasis als Hauptbegrenzung. Die Modellinferenz ist mit 85\,ms auf der fusionierten Punktwolke 10-Hz-fähig; die Echtzeitfähigkeit des vollständigen Online-Fusionssystems ist dagegen nicht nachgewiesen, da der vorgelagerte Merge bislang nur als Offline-Pipeline mit 2{,}25\,s pro Frame existiert.

\bigskip

Die Fusion ist nicht erforderlich, um das bewegte Auto zu erkennen. In der experimentweisen Auswertung ist \texttt{os0} auf allen drei Autoaufnahmen sogar die stärkste Ansicht (AP\textsubscript{30} zwischen 0{,}92 und 0{,}96); in der gepoolten Auswertung über alle Testvorhersagen liegen die drei Ansichten mit 0{,}891, 0{,}885 und 0{,}832 dicht beieinander. Der frühere starke \texttt{os1}-Einbruch war ein Artefakt der zeitlich und räumlich gekoppelten Testaufteilung. Dieses Ergebnis unterstreicht, wie stark wissenschaftliche Schlussfolgerungen von einem passenden Evaluationsdesign abhängen.

Die aussagekräftigste Gesamtformulierung lautet:

\begin{quote}
    Die Multisensorfusion liefert die beste ausgewogene Gesamtleistung und erreicht bei Personen und Fahrrädern in beiden Auswertungsvarianten die präziseste Lokalisierung, ist aber nicht in jeder Klasse und bei jeder Metrik die beste Ansicht.
\end{quote}

Für einen späteren Einsatz sind vor allem ein finaler Refit des ausgewählten Modells, ein gemessener Online-Merge, validierungsbasierte Betriebsschwellen und zeitliches Tracking erforderlich.
